% Complete technical appendix. Input after \appendix in the main manuscript.
% Requires amsmath, amssymb, amsthm and theorem/proposition/lemma/corollary environments.

\section{Protocol, estimands, and complete proofs}
\label{app:theory}

This appendix gives complete proofs for the proposed protocol. The
concentration and betting arguments are specializations of established
anytime-valid inference, rather than new concentration inequalities
\citep{howard2021confidence,waudbysmith2024betting}. Sequential win-statistic
methods already exist \citep{zhang2024sequential,bergemann2026group}, as do
asymptotic anytime-valid methods for U-statistics
\citep{cai2026ustatistics} and confidence-sequence-based deployment gates
\citep{karampatziakis2021offpolicy}. The purpose of the results below is to
make the evaluation design, decision claim, and its limitations explicit.

\subsection{A measurable hierarchy with operational ties}

Let $Y\in\mathcal Y$ denote the complete outcome record of an episode under
a frozen evaluation protocol. It includes final outcomes, selected
trajectory summaries, and resources measured through a fixed horizon.
Let $d_k:\mathcal Y^2\to\{-1,0,1\}$, $k=1,\ldots,K$, be measurable,
antisymmetric tier comparators:
$d_k(y,y')=-d_k(y',y)$. A value of zero means a prespecified operational
tie. For example, for a larger-is-better scalar endpoint $g_k$ and an
absolute tolerance $\eta_k\geq0$,
\[
 d_k(y,y')=\mathbf1\{g_k(y)-g_k(y')>\eta_k\}
           -\mathbf1\{g_k(y')-g_k(y)>\eta_k\}.
\]
Resource tiers can be set to zero unless a symmetric eligibility condition
is satisfied, such as both episodes achieving verified success. Such a
rule is part of the estimand. It must not be chosen after inspecting the
systems' comparative results. Define
\begin{align}
 h_k(y,y')&=d_k(y,y')\prod_{\ell<k}\mathbf1\{d_\ell(y,y')=0\},
 & h(y,y')&=\sum_{k=1}^K h_k(y,y'). \label{eq:hierarchy_kernel}
\end{align}
At most one term is nonzero; hence $h\in\{-1,0,1\}$ and $h$ is
antisymmetric. The first decisive tier determines the sign. A tolerance
does not assert statistical equivalence, and absence of a significant
difference is not used to define an episode-level tie.

For any specified joint comparison law $P$ on $(Y^A,Y^B)$, put
\[
 p_W=P(h=1),\quad p_L=P(h=-1),\quad p_T=P(h=0),\quad
 \theta=E_P h=p_W-p_L.
\]
The net benefit $\theta$ is the primary bounded target. Tier contributions
are $\theta_k=E_P h_k$, so that $\theta=\sum_k\theta_k$. The corresponding
decisive masses $P(|h_k|=1)$ show how often each tier determined a
comparison; these are not population-level component effects.

\begin{proposition}[Relations among win summaries]\label{prop:win_relations}
Let $d=p_W+p_L$. When denominators are positive,
\[
 \mathrm{WR}=\frac{d+\theta}{d-\theta},\qquad
 \mathrm{WO}=\frac{p_W+p_T/2}{p_L+p_T/2}
             =\frac{1+\theta}{1-\theta}.
\]
In particular, $\mathrm{WR}>1$ is equivalent to $\theta>0$ when $p_L>0$.
The net benefit alone determines win odds, but does not determine the win
ratio.
\end{proposition}
\begin{proof}
Adding and subtracting $d=p_W+p_L$ and $\theta=p_W-p_L$ gives
$p_W=(d+\theta)/2$ and $p_L=(d-\theta)/2$. Also,
$p_W+p_T/2=(1+\theta)/2$ and $p_L+p_T/2=(1-\theta)/2$.
The identities and equivalence follow. Different $d\in[|\theta|,1]$
give different win ratios at a fixed nonzero $\theta$.
\end{proof}

\subsection{Pairwise priority does not imply a population guardrail}

\begin{proposition}[Arbitrarily favorable preference can hide harm]
\label{prop:compensation}
Consider verified success followed by cost, with lower cost preferred
only when both systems succeed. For every $\epsilon\in(0,1)$ there is a
comparison distribution with success-rate difference
$\Delta_S=P(S^A=1)-P(S^B=1)=-\epsilon$ and net benefit
$\theta=1-2\epsilon$. Consequently, for any noninferiority margin
$\delta\in[0,1/2)$, positive net benefit need not imply
$\Delta_S>-\delta$. Moreover, net benefit can approach one while the
success-rate difference remains strictly negative.
\end{proposition}
\begin{proof}
Let $B$ always succeed. Let $A$ fail with probability $\epsilon$ and
otherwise succeed with lower cost than $B$. On failure, $A$ loses at the
success tier. On success, $A$ wins at the cost tier. Thus $p_W=1-\epsilon$,
$p_L=\epsilon$, and $\Delta_S=-\epsilon$. For the margin claim choose
$\delta<\epsilon<1/2$. Taking $\epsilon\downarrow0$ proves the last
claim. This construction works for same-task or independent-draw
comparisons because $B$'s outcome is constant.
\end{proof}

The proposition does not say that cost should be ignored. It shows that
``success takes priority within a pair'' and ``population success cannot
decline beyond a margin'' are different requirements. We therefore use
component scores $g_j:\mathcal Y\to[0,1]$, oriented so larger values are
better, and separate component differences $\Delta_j$. A safety score can
be one minus an adverse-event indicator. An adverse-event rate that is
already unacceptable for the reference also requires an \emph{absolute}
rate gate: noninferiority alone does not certify an acceptable absolute
rate. The bounded-score results below also apply to an absolute score for
$A$, by replacing a difference kernel with $g_j(y^A)$.

\subsection{What an online A/B experiment identifies}

A live-episode counterfactual, a same-task clone comparison, and a
two-arrival online comparison are three distinct targets. Let
$Y^A,Y^B$ first denote the potential outcomes of a single live episode
and define $\theta_{\mathrm{episode}}=E h(Y^A,Y^B)$. This target depends
on the joint potential-outcome law, including latent episode features.
For the clone design, let $Y^a(x,\xi)$ be a complete
episode under system $a$, task context $x$, and evaluation randomness
$\xi$. The offline independent-replicate target is
\begin{equation}
 \theta_{\mathrm{task}}
 =E_X E_{\xi_A,\xi_B\mid X}
       h\{Y^A(X,\xi_A),Y^B(X,\xi_B)\}, \label{eq:task_target}
\end{equation}
where the two replicate streams are conditionally independent. A
common-randomness coupling defines a different target unless equality of
these expectations is established. The law of $X$, evaluation horizon,
and replication convention must all be reported.

For a stratum $s$, an online independent-draw target is
\begin{equation}
 \theta_{\mathrm{cross},s}
 = E\,h(Y_s^A,\widetilde Y_s^B),\qquad
 Y_s^A\perp\widetilde Y_s^B, \label{eq:cross_target}
\end{equation}
where each episode is drawn from that stratum's workload. Even exact
knowledge of both marginal outcome distributions generally does not
identify $\theta_{\mathrm{episode}}$. By contrast, if the arm-specific
conditional distributions given an observed task context $X$ are
identified under overlap, their product defines the identifiable
functional in \eqref{eq:task_target}. Individual-counterfactual,
same-covariate independent-copy, and population independent-copy win
targets, including the first target's nonidentifiability, are explicitly
distinguished by \citet{even2025rethinking}. Our construction below
illustrates that established obstruction for an agent episode.

\begin{proposition}[Episode-counterfactual preference is not identified]
\label{prop:nonidentification}
There exist two full-data laws with identical distributions of all
observed data from an independently randomized, one-system-per-task
A/B experiment, but opposite $\theta_{\mathrm{episode}}$. Both laws have
independent-draw net benefit zero.
\end{proposition}
\begin{proof}
Let the observed context $X$ be constant, let the latent episode feature
$U$ be uniform on $\{0,1,2\}$, and use the single-tier comparator
$h(a,b)=\operatorname{sign}(a-b)$. Under the first law set
$Y^B=U$ and $Y^A=(U+1)\bmod3$. Under the second set
$Y^B=U$ and $Y^A=(U-1)\bmod3$. Under either law, each system's marginal
outcome is uniform on $\{0,1,2\}$. A randomized assignment independent of
potential outcomes therefore gives the same distribution for
$(\text{assignment},\text{observed outcome})$; independent tasks give the
same law at every sample size. In the first law, $A$ wins for $U=0,1$
and loses for $U=2$, giving $\theta_{\mathrm{episode}}=1/3$. In the second,
$A$ wins for $U=0$ and loses for $U=1,2$, giving $-1/3$.
For independent uniform draws the probability of either strict ordering
is $1/3$, so both $\theta_{\mathrm{cross}}$ and the independent-replicate
$\theta_{\mathrm{task}}$ given the constant observed context are zero
under both laws.
\end{proof}

This is an identification obstruction, not a confidence-interval issue.
Additional observations of the same randomized design cannot resolve
$\theta_{\mathrm{episode}}$.
A binary single-tier net benefit is a special exception because it
equals the difference of the two means; the proposition uses three
ordered values. Running both systems in resettable task environments can
identify the explicitly chosen clone/run coupling, subject to valid
measurement and independent task sampling. Shadow runs do not reproduce
live user reactions or downstream interference automatically.

\subsection{A randomized disjoint-pair design}

Index nonoverlapping pairs of enrolled episodes by $i$. The pair has two
positions, $r=1,2$, and potential outcomes $Y_{ir}^A,Y_{ir}^B$.
Both positions are fixed before randomizing the orientation. All
potential outcomes here refer to the fixed system versions, endpoint
definitions, horizon, and resource caps. There is no interference between
the two positions. Let $R_i=1$ assign $A$ to position 1 and $B$ to
position 2; $R_i=0$ reverses these assignments. Let
\[
 P(R_i=1\mid\mathcal H_i,\{Y_{ir}^a\}_{r,a})=q_i,
 \qquad 0<\varepsilon\leq q_i\leq1-\varepsilon,
\]
where $\mathcal H_i$ contains the information used to form and randomize
the pair, and $q_i$ is known and $\mathcal H_i$-measurable. This condition
is satisfied by a fresh randomization coin with its recorded probability.
It excludes assignments driven by unrecorded current outcomes.

For any measurable kernel $\phi:\mathcal Y^2\to[-1,1]$, define observed
outcomes $O_{i1},O_{i2}$ and the score
\begin{equation}
 Z_i^\phi=
 \frac{R_i}{2q_i}\phi(O_{i1},O_{i2})+
 \frac{1-R_i}{2(1-q_i)}\phi(O_{i2},O_{i1}).
 \label{eq:pair_ht}
\end{equation}
The arguments are always oriented as $A$ then $B$. The within-pair
symmetrized potential-outcome quantity is
\begin{equation}
 m_i^\phi=\frac12\{
 \phi(Y_{i1}^A,Y_{i2}^B)+\phi(Y_{i2}^A,Y_{i1}^B)\}.
 \label{eq:finite_pair_target}
\end{equation}

\begin{theorem}[Identification by pair randomization]\label{thm:pair_id}
Under the stated randomization and consistency conditions,
\[
 E(Z_i^\phi\mid\mathcal H_i,\{Y_{ir}^a\}_{r,a})=m_i^\phi.
\]
For a filtration $\mathcal F_{i-1}$ describing prior pairs and satisfying
$\mathcal F_{i-1}\subseteq\mathcal H_i$, it follows that
$E(Z_i^\phi\mid\mathcal F_{i-1})
=E(m_i^\phi\mid\mathcal F_{i-1})=:\mu_i^\phi$.
The score is bounded by $|Z_i^\phi|\leq1/(2\varepsilon)$.
With balanced orientation $q_i=1/2$, it is simply the observed oriented
kernel, and $|Z_i^\phi|\leq1$.
\end{theorem}
\begin{proof}
Under $R_i=1$, consistency gives
$(O_{i1},O_{i2})=(Y_{i1}^A,Y_{i2}^B)$. Its contribution to the
conditional expectation is
$q_i\phi(Y_{i1}^A,Y_{i2}^B)/(2q_i)$.
Under $R_i=0$, the contribution is
$(1-q_i)\phi(Y_{i2}^A,Y_{i1}^B)/(2(1-q_i))$.
Their sum is \eqref{eq:finite_pair_target}. The conditional-mean identity
follows by iterated expectation. Exactly one term in
\eqref{eq:pair_ht} is nonzero, yielding the stated bound.
\end{proof}

Conditional exchangeability of the two positions is \emph{not} needed
for the symmetrized target in \eqref{eq:finite_pair_target}. If,
additionally, the two potential-outcome vectors are conditionally
independent draws from the same stratum-specific population, each term's
conditional mean equals \eqref{eq:cross_target}. Then balanced pairs
estimate the independent-draw stratum preference. Pair formation by
baseline stratum reduces cross-context comparisons but does not create
same-task observations.

For a population component set
$\phi_j(y,y')=g_j(y)-g_j(y')$. Algebra gives the particularly useful
identity
\begin{equation}
 m_i^{\phi_j}=
 \frac12\sum_{r=1}^2\{g_j(Y_{ir}^A)-g_j(Y_{ir}^B)\}.
 \label{eq:component_pair_effect}
\end{equation}
Unlike a general hierarchy kernel, the component contrast is the
average same-episode component effect for the pair. Thus the design
supports population component gates and a cross-episode preference
without incorrectly claiming a same-task preference effect.

If $q_i\ne1/2$, omitting the inverse-probability factors generally
targets $q_i\phi(Y_{i1}^A,Y_{i2}^B)
+(1-q_i)\phi(Y_{i2}^A,Y_{i1}^B)$ instead. For example, if these two
potential comparison signs are $1$ and $-1$ and $q_i=3/4$, the
symmetrized target is zero but the unweighted mean is $1/2$.
The extension permits adaptive \emph{orientation} probabilities; it
still allocates one episode to each system per pair and is not a general
bandit-router analysis.

\subsection{Anytime confidence sequences, including drift}

The next result applies to any score sequence $Z_i$, including each
hierarchy, tier, probability, or component score above. Let
$(\mathcal F_i)$ be a filtration to which the scores are adapted.
Assume known $\mathcal F_{i-1}$-measurable bounds
$a_i\leq Z_i\leq b_i$, set $c_i=b_i-a_i$, and write
\[
 \mu_i=E(Z_i\mid\mathcal F_{i-1}),\quad
 \overline\mu_n=\frac1n\sum_{i=1}^n\mu_i,\quad
 \widehat\mu_n=\frac1n\sum_{i=1}^n Z_i,\quad
 V_n=\frac14\sum_{i=1}^n c_i^2.
\]
The predictable bounds must hold before observing $Z_i$. If $q_i$ depends
on current contexts outside $\mathcal F_{i-1}$, the universal
$[-1/(2\varepsilon),1/(2\varepsilon)]$ bound is valid; substituting a
post-context narrower bound requires a filtration and centering
argument that preserve predictability.

\begin{theorem}[Normal-mixture confidence sequence]\label{thm:normal_cs}
Fix $\rho>0$ and $\alpha\in(0,1)$ before examining this score stream.
Define
\begin{equation}
 B_\alpha(v)=
 \sqrt{(v+\rho)\log\!\left(\frac{v+\rho}{\rho\alpha^2}\right)},
 \qquad
 C_n=\left[\widehat\mu_n-\frac{B_\alpha(V_n)}n,
           \widehat\mu_n+\frac{B_\alpha(V_n)}n\right].
 \label{eq:normal_cs}
\end{equation}
Then
\[
 P\{\overline\mu_n\in C_n\text{ for every }n\geq1\}\geq1-\alpha.
\]
One may intersect $C_n$ with a known parameter range, such as $[-1,1]$.
For balanced pair scores in $[-1,1]$, $V_n=n$.
\end{theorem}
\begin{proof}
Let $D_i=Z_i-\mu_i$ and $S_n=\sum_{i\leq n}D_i$.
The conditional Hoeffding lemma gives, for every fixed
$\lambda\in\mathbb R$,
\[
 E\{e^{\lambda D_i}\mid\mathcal F_{i-1}\}
 \leq e^{\lambda^2c_i^2/8}.
\]
Therefore $M_n(\lambda)=\exp\{\lambda S_n-\lambda^2V_n/2\}$ is a
nonnegative supermartingale starting at one. Mix over
$\lambda\sim N(0,\rho^{-1})$. Conditional Tonelli's theorem preserves
the supermartingale property. Completing the square in the Gaussian
integral gives
\[
 M_n=\sqrt{\frac\rho{V_n+\rho}}
      \exp\!\left\{\frac{S_n^2}{2(V_n+\rho)}\right\}.
\]
Ville's inequality implies $P(\exists n:M_n\geq1/\alpha)\leq\alpha$.
If $|S_n|\geq B_\alpha(V_n)$, substitution gives
$M_n\geq1/\alpha$. Hence, outside an event of probability at most
$\alpha$, $|\widehat\mu_n-\overline\mu_n|<B_\alpha(V_n)/n$ for every
$n$. This proves the claimed coverage. Intersecting with a set known
to contain the target cannot remove it on the coverage event.
\end{proof}

No identical-distribution assumption appears in the theorem. Under
drift, however, the target is the \emph{running average conditional
effect} $\overline\mu_n$, not the effect at the latest arrival or a
future workload. If a predictable policy changes the stratum mix,
hierarchy, or candidate version, that change is reflected in the sequence
$\mu_i$. Predictability preserves the concentration argument but does
not preserve a frozen deployment estimand. A new version or preference
rule should ordinarily start a separately specified experiment.

The normal-mixture width scales as $O(\sqrt{\log n/n})$ for bounded
constant ranges and fixed $\rho$. It is convenient and explicit, but
is not an iterated-logarithm-optimal boundary. Tuning $\rho$ after
examining this stream invalidates this particular guarantee unless that
adaptation is accounted for.

\subsection{Betting tests for stationary deployment claims}

For stationary claims, a direct test can be more efficient than the
range-based confidence sequence. Let $Z_i\in[-1,1]$ and test
$H_0:E(Z_i\mid\mathcal F_{i-1})\leq c$ for every $i$, where
$c\in(-1,1)$ is fixed. Select $J$ nonnegative stakes satisfying
$0\leq\lambda_j<1/(1+c)$ and mixture weights
$w_j>0$, $\sum_jw_j=1$, before examining the stream. Define
\begin{equation}
 E_n(c)=\sum_{j=1}^J w_j
             \prod_{i=1}^n\{1+\lambda_j(Z_i-c)\}.
 \label{eq:betting_e}
\end{equation}
These are standard bounded-mean betting tests
\citep{waudbysmith2024betting}; sequential two-sample betting also has an
established literature \citep{shekhar2024betting,podkopaev2023predictive}.

\begin{theorem}[Validity and the finite-grid power condition]
\label{thm:betting}
Under $H_0$, $E_n(c)$ is a nonnegative supermartingale starting at one,
and
$P\{\sup_n E_n(c)\geq1/\alpha\}\leq\alpha$.
If the scores are i.i.d. under an alternative and at least one selected
stake satisfies
$E\log\{1+\lambda_j(Z_1-c)\}>0$, then $E_n(c)\to\infty$ almost surely.
A fixed finite grid need not satisfy this condition for every alternative
with $E Z_1>c$.
\end{theorem}
\begin{proof}
Each factor is positive because
$1+\lambda_j(Z_i-c)\geq1-\lambda_j(1+c)>0$.
Its conditional expectation is
$1+\lambda_j\{E(Z_i\mid\mathcal F_{i-1})-c\}\leq1$.
Each product is therefore a nonnegative supermartingale; their fixed
convex mixture is also one. Ville's inequality proves the error bound.
For a selected stake with positive expected log factor, the factors are
bounded away from zero and infinity, so their logarithms are integrable.
The strong law implies that the corresponding log product divided by
$n$ tends to its strictly positive expectation. That product and hence
its positively weighted contribution to $E_n(c)$ diverge.
For the final claim take $c=0$ and a nondegenerate symmetric
$Z_1\in\{-1,1\}$. Every fixed $\lambda\in(0,1)$ has negative mean
log factor, since it equals $\tfrac12\log(1-\lambda^2)$.
Continuity in the success probability implies that all finitely many
positive grid stakes still have negative mean log factor for a sufficiently
small positive mean. Zero stakes remain constant and cannot supply
divergence.
\end{proof}

A countable positive-weight mixture with stakes tending to zero repairs
the latter consistency limitation: for $EZ_1-c>0$, the derivative at
zero of $E\log\{1+\lambda(Z_1-c)\}$ is $EZ_1-c>0$, so all sufficiently
small positive stakes have positive mean log factor. A reported finite
grid retains exact error control, but its stake choices and empirical
power must be reported without a universal consistency claim.
The null in Theorem~\ref{thm:betting} is a conditional-mean null at
every step. Merely requiring a nonpositive current cumulative mean under
arbitrary drift does not imply the product's supermartingale property.
For that running-average target use Theorem~\ref{thm:normal_cs}.

\subsection{Conjunctive deployment and the multiplicity distinction}

Index the required criteria by $j=0,\ldots,m-1$. Criterion zero is
hierarchical net benefit above a prespecified minimum $c_0\geq0$.
Other criteria can be component noninferiority,
$\Delta_j>c_j=-\delta_j$, with $\delta_j\geq0$. A deployment claim is
the conjunction that \emph{every} criterion holds. Use the corresponding
bounded scores from \eqref{eq:pair_ht}. Let $L_{jn}(\alpha_j)$ be the
lower endpoint in \eqref{eq:normal_cs}. The rule is
\begin{equation}
 \tau=\inf\{n\geq1:L_{jn}(\alpha_j)>c_j
                     \text{ for all }j\}. \label{eq:deploy_rule}
\end{equation}
The reported stationary betting implementation replaces each lower-bound
condition with $E_{jn}(c_j)\geq1/\alpha_j$: all criteria must pass at the
same monitored look. An optional retained-evidence variant instead uses
$\sup_{t\leq n}E_{jt}(c_j)\geq1/\alpha_j$. Retention is justified only
for the fixed stationary claims; its stopping times and power differ
from those of the reported current-look implementation.

\begin{theorem}[Stationary conjunction needs no alpha split]
\label{thm:iut}
Suppose each score has fixed conditional mean
$E(Z_{ij}\mid\mathcal F_{i-1})=\mu_j$ and at least one requirement is
false, namely $\mu_{j_*}\leq c_{j_*}$. Set $\alpha_j=\alpha$ for every
criterion. The confidence-sequence rule and both stationary betting variants have
$P(\tau<\infty)\leq\alpha$, regardless of dependence among criteria.
\end{theorem}
\begin{proof}
For the confidence-sequence rule, deployment requires
$L_{j_*n}(\alpha)>c_{j_*}\geq\mu_{j_*}$ at some $n$. Its probability
is at most $\alpha$ by Theorem~\ref{thm:normal_cs}.
For betting, deployment requires the $j_*$ process to cross
$1/\alpha$; its conditional-mean null holds, so
Theorem~\ref{thm:betting} bounds this probability by $\alpha$.
The index $j_*$ can be chosen for each fixed data-generating law. No
independence among criteria or union bound over them is needed.
\end{proof}

This is an intersection--union test for a \emph{single conjunction},
not simultaneous confidence coverage for all component claims. It does
not authorize choosing among multiple candidates, hierarchies, success
routes, or deployment experiments at level $\alpha$ each. For example,
independent repeated experiments each allowed a level-$\alpha$ false
deployment can have probability $1-(1-\alpha)^M$ of at least one false
deployment over $M$ experiments. An experiment-wise allocation with
$\sum_e\alpha_e\leq\alpha_{\mathrm{program}}$ gives a simple overall
bound by the union inequality.

\begin{theorem}[Drifting running-average claims require simultaneous control]
\label{thm:drift_gate}
Let $\overline\mu_{jn}=n^{-1}\sum_{i\leq n}
E(Z_{ij}\mid\mathcal F_{i-1})$, with arbitrary allowed drift.
If $\sum_j\alpha_j\leq\alpha$, the confidence-sequence rule satisfies
\[
 P\{\exists n:\ \text{the rule deploys at }n
          \text{ and some }\overline\mu_{jn}\leq c_j\}\leq\alpha.
\]
Alternatively, level $\alpha$ per criterion remains sufficient if
there is one fixed index $j_*$ such that
$\overline\mu_{j_*n}\leq c_{j_*}$ almost surely for every $n$.
\end{theorem}
\begin{proof}
The union bound and Theorem~\ref{thm:normal_cs} give probability at least
$1-\sum_j\alpha_j$ that all lower bounds are valid for all indices and
times. On this event deployment implies
$\overline\mu_{jn}\geq L_{jn}>c_j$ for every $j$, so the stated error
event cannot occur. Under the fixed-index condition the error event
is contained in failure of the one $j_*$ lower confidence sequence,
which has probability at most $\alpha$.
\end{proof}

Without a fixed violating component, a time-varying conjunction cannot
use the stationary proof: the relevant failing component can change
with $n$. Nor does a favorable historical average establish that a system
will be safe after a workload shift. A finite-sample statement about the
future needs a stability or transport assumption beyond the results here.

\subsection{Delay, strata, and versions}

\begin{proposition}[Displaying completed enrollment prefixes]
\label{prop:delay}
Suppose the full episode scores in enrollment-pair order satisfy
Theorem~\ref{thm:normal_cs}. At calendar time $t$, let $N(t)$ be any
possibly data-dependent integer such that the first $N(t)$ pairs have
their frozen complete endpoints available. Display $C_{N(t)}$, with no
display at index zero. The same simultaneous coverage guarantee holds
over all calendar display times.
\end{proposition}
\begin{proof}
On the event that the confidence sequence covers at every integer
$n\geq1$, it covers at every selected integer $N(t)\geq1$, simultaneously
over $t$. Thus every calendar-time coverage failure is contained in the
original failure event. No stopping-time property of $N(t)$ is needed
for this pathwise inclusion.
\end{proof}

The premise is substantive: the score's conditional-mean and bound
assumptions must hold in the filtration used for enrollment-order
inference. Shared load, interference, and adaptive collection cannot
be dismissed by invoking the delay proposition. A practical protocol
fixes a horizon and treats failure to achieve a required outcome by that
horizon as an endpoint, not as an observation to discard. It waits for
each pair to mature and analyzes completed enrollment prefixes. Naively
sorting by completion time or selecting only completed successes can
change the comparison population. If endpoints remain genuinely missing
at the fixed horizon, this analysis requires a prespecified observed
missingness outcome or a separate justified missing-data model; it does
not impute validity to missing traces.

For a prespecified set of $S$ strata with stationary parameters $\mu_s$
and a fixed target mix $w_s\geq0$, $\sum_sw_s=1$, construct within-stratum
confidence sequences $[L_{s,n_s},U_{s,n_s}]$ at levels $\alpha_s$ with
$\sum_s\alpha_s\leq\alpha$. Set $[L_{s,0},U_{s,0}]=[-1,1]$ before a
stratum is sampled. Then
\begin{equation}
 \left[\sum_s w_sL_{s,n_s},\ \sum_s w_sU_{s,n_s}\right]
 \label{eq:stratified_cs}
\end{equation}
covers $\sum_sw_s\mu_s$ for every vector of sample counts with probability
at least $1-\alpha$. Indeed, simultaneous coverage of all stratum
sequences implies the result by nonnegative weighting. Adaptive
allocation of sample counts is allowed provided each within-stratum
sequence retains its stated assumptions. Data-dependent weights selected
to favor the candidate do not preserve a prespecified workload target.
If both strata and gates require simultaneous drifting claims, allocate
error over both indices.

\subsection{Offline repeated runs: the task is the independent unit}

Let $X_i$, $i=1,\ldots,n$, be i.i.d. task clusters. The complete records
consisting of task information, replicate counts, and all episode
randomness are independent across $i$. Conditional on $X_i$,
draw independent replicate sets
$Y_{ir}^A$, $r=1,\ldots,R_{iA}$, and
$Y_{is}^B$, $s=1,\ldots,R_{iB}$, with the stated system-specific laws.
The positive replicate counts may be fixed or determined by task
information before observing its outcomes. Define
\begin{equation}
 T_i=\frac1{R_{iA}R_{iB}}
       \sum_{r=1}^{R_{iA}}\sum_{s=1}^{R_{iB}}
               h(Y_{ir}^A,Y_{is}^B),\qquad
 \widehat\theta_n=\frac1n\sum_{i=1}^nT_i.
 \label{eq:offline_est}
\end{equation}

\begin{theorem}[Task-level estimation and inference]\label{thm:offline}
Under this design, $E(T_i)=\theta_{\mathrm{task}}$,
$T_i\in[-1,1]$, and $\widehat\theta_n$ is unbiased. If the entire
task-and-replication records are also identically distributed across tasks,
then
\[
 \operatorname{Var}(\widehat\theta_n)=\frac{\operatorname{Var}(T_1)}n.
\]
The normal-mixture confidence sequence with scores $T_i$ and $V_n=n$
is finite-sample valid in task order. Under the preceding i.i.d.
complete-record condition, if
$\sigma^2=\operatorname{Var}(T_1)>0$, then
\[
 \frac{\sqrt n(\widehat\theta_n-\theta_{\mathrm{task}})}
 {\{(n-1)^{-1}\sum_i(T_i-\widehat\theta_n)^2\}^{1/2}}
 \ \Longrightarrow\ N(0,1).
\]
\end{theorem}
\begin{proof}
Each replicate comparison has the conditional expectation in
\eqref{eq:task_target}; averaging and iterating expectation proves
unbiasedness. Boundedness follows because $T_i$ is an average of
numbers in $[-1,1]$. Independence across tasks makes the variance of
the mean equal to the sum of the task variances divided by $n^2$,
which simplifies as stated under identical distribution. The bounded
task scores meet Theorem~\ref{thm:normal_cs}. Finally, the ordinary
i.i.d. central limit theorem applies because $T_i$ is bounded. The
sample variance converges to $\sigma^2$ by the strong law, and
Slutsky's theorem gives the studentized limit.
\end{proof}

Replicate pairs reuse episode outcomes and are dependent. Their number
$\sum_iR_{iA}R_{iB}$ is not an independent sample size. If tasks share a
repository, user, or scenario with residual dependence, independent
clusters at that higher level must replace individual tasks. If the
benchmark is a fixed list rather than a sample of tasks, generalization
to an external workload does not follow from a task bootstrap; one must
state a fixed-list target or justify task sampling. Repeated runs do not
repair that generalization gap.

For $\mathrm{WR}=p_W/p_L$, use two task-level scores formed by replacing
$h$ with $\mathbf1\{h=1\}$ and $\mathbf1\{h=-1\}$. Simultaneous
intervals $[l_W,u_W]$ and $[l_L,u_L]$ with $l_L>0$ imply
$\mathrm{WR}\in[l_W/u_L,u_W/l_L]$. If the lower loss bound is zero,
there is no finite upper ratio bound from these intervals. In contrast,
any confidence interval for $\theta$ transforms monotonically to win
odds on $(-1,1)$ via Proposition~\ref{prop:win_relations}.

\begin{lemma}[Shared seeds: the diagonal changes the target]
\label{lem:shared_seed}
For a fixed task $X=x$, suppose the paired seed records
$W_r=(Y_r^A,Y_r^B)$, $r=1,\ldots,R$, are i.i.d. conditional on $x$,
but the two outcomes within $W_r$ may be dependent. Let
\[
 \theta_{\mathrm{ind}}(x)
 =E\{h(Y_1^A,Y_2^B)\mid x\},\qquad
 \theta_{\mathrm{coupled}}(x)
 =E\{h(Y_1^A,Y_1^B)\mid x\}.
\]
For $R\geq2$, the off-diagonal statistic
\[
 U_R(x)=\frac1{R(R-1)}\sum_{r\ne s}h(Y_r^A,Y_s^B)
\]
is conditionally unbiased for $\theta_{\mathrm{ind}}(x)$.
The all-pairs statistic, including diagonal comparisons, satisfies
\[
 E\!\left\{\frac1{R^2}\sum_{r,s}h(Y_r^A,Y_s^B)\,\middle|\,x\right\}
 =\left(1-\frac1R\right)\theta_{\mathrm{ind}}(x)
       +\frac1R\theta_{\mathrm{coupled}}(x).
\]
Its absolute bias for $\theta_{\mathrm{ind}}(x)$ is at most $2/R$.
\end{lemma}
\begin{proof}
For $r\ne s$, conditional independence of seed records gives
$E\{h(Y_r^A,Y_s^B)\mid x\}=\theta_{\mathrm{ind}}(x)$.
For $r=s$, the expectation is
$\theta_{\mathrm{coupled}}(x)$. There are $R(R-1)$ off-diagonal
terms and $R$ diagonal terms. Averaging proves both identities. Since
both expectations lie in $[-1,1]$, their difference divided by $R$
has absolute value at most $2/R$.
\end{proof}

This is the usual finite-sample distinction between a U-statistic and
its diagonal-including V-statistic, applied to paired stochastic agent
runs. Averaging $U_R(X_i)$ over independent tasks retains the bounded
task-level inference in Theorem~\ref{thm:offline}. The
$R(R-1)$ comparisons still are not independent observations.
For $R=4$, the all-pairs expectation contains one-quarter coupled-seed
preference. With only one seed, no off-diagonal estimate is available.
Removing diagonal pairs repairs within-seed coupling only under
independence of the distinct seed records; shared infrastructure or
other cross-seed dependence needs its own treatment.

\subsection{Sampling cost of requiring every gate}

\begin{proposition}[A finite-time tail bound for deployment delay]
\label{prop:delay_bound}
Suppose stationary criterion $j$ has conditional mean
$\mu_j=c_j+\Delta_j$ with $\Delta_j>0$ and constant score-range width
$d_j$. Write $r_{jn}=B_{\alpha_j}(nd_j^2/4)/n$ for its confidence
radius. For any deterministic $n$ satisfying $r_{jn}<\Delta_j$ for all
$j$, the confidence-sequence deployment time in
\eqref{eq:deploy_rule} satisfies
\[
 P(\tau>n)\leq\sum_{j=0}^{m-1}
       \exp\!\left\{-\frac{2n(\Delta_j-r_{jn})^2}{d_j^2}\right\}.
\]
In particular, for fixed positive gaps and bounded widths, deployment
occurs almost surely as data continue to accrue.
\end{proposition}
\begin{proof}
If every $L_{jn}>c_j$, the rule deploys by $n$. For a given criterion,
$L_{jn}\leq c_j$ implies
$\widehat\mu_{jn}-\mu_j\leq-(\Delta_j-r_{jn})$.
The conditional Hoeffding lemma followed by the Chernoff bound gives
probability at most
$\exp\{-2n(\Delta_j-r_{jn})^2/d_j^2\}$ for that event. The union
bound proves the result without requiring independence between criteria.
Because $r_{jn}\to0$, the right side tends to zero. The events
$\{\tau>n\}$ decrease to $\{\tau=\infty\}$, so continuity of
probability proves almost sure finite deployment.
\end{proof}

This bound is deliberately conservative. It makes explicit that a
criterion close to its margin can dominate the sample requirement even
when hierarchical preference is overwhelmingly favorable.

\begin{proposition}[A guardrail has an unavoidable information cost]
\label{prop:guardrail_lower_bound}
Consider a submodel in which each pair reveals one $A$ failure
indicator $F_i\sim\mathrm{Bernoulli}(p)$ independently, $B$ never
fails, and all other recorded information has the same conditional law
under the hypotheses and gives no additional information about $p$.
Fix a permitted failure rate $\delta\in(0,1)$ and an acceptable
alternative $p_1\in(0,\delta)$. Any sequential procedure with almost
surely finite decision time $T$ under both laws, $E_{p_1}T<\infty$,
and deployment event $D$ satisfying
$P_{\delta}(D)\leq\alpha<1-\beta\leq P_{p_1}(D)$ must obey
\[
 E_{p_1}T\ \geq\
 \frac{\operatorname{kl}(1-\beta,\alpha)}
      {\operatorname{kl}(p_1,\delta)},
 \qquad
 \operatorname{kl}(u,v)=u\log\frac uv+(1-u)\log\frac{1-u}{1-v}.
\]
\end{proposition}
\begin{proof}
Let $\ell_i$ be the log likelihood ratio of one failure observation
under $p_1$ versus $\delta$. The stopped transcript likelihood ratio
is $\exp(\sum_{i=1}^T\ell_i)$, since the procedure's randomization
mechanism is the same under both laws. Its relative entropy under
$p_1$ is
\[
 E_{p_1}\sum_{i=1}^T\ell_i
 =\sum_{i\geq1} E_{p_1}[\mathbf1\{T\geq i\}\ell_i]
 =E_{p_1}T\,\operatorname{kl}(p_1,\delta).
\]
The interchange is valid because $\ell_i$ is bounded for interior
$p_1,\delta$ and $E_{p_1}T<\infty$. The final equality uses the
independence of $F_i$ from the pre-$i$ stopping decision.
Mapping the transcript to $\mathbf1_D$ cannot increase relative
entropy (the log-sum inequality), so this quantity is at least
$\operatorname{kl}\{P_{p_1}(D),P_\delta(D)\}$.
For $u>v$, binary relative entropy increases with $u$ and decreases
with $v$, as follows by differentiating its displayed formula.
The stated probability bounds therefore give a lower bound of
$\operatorname{kl}(1-\beta,\alpha)$, proving the claim.
\end{proof}

For $p_1=\delta-\gamma$ with $\gamma\to0$ and fixed interior $\delta$,
Taylor expansion yields
$\operatorname{kl}(\delta-\gamma,\delta)
=\gamma^2/\{2\delta(1-\delta)\}+o(\gamma^2)$.
Thus the required expected sample size can scale as the inverse squared
distance to the guardrail. The result is a standard information-bound
argument applied to this deployment subproblem, not a claim of a new
general sequential lower bound.

\subsection{Sensitivity of a preference claim to its comparator}

\begin{proposition}[Comparator disagreement bound]
\label{prop:sensitivity}
For any two prespecified measurable comparators $h,g\in[-1,1]$ on the
same comparison population,
\[
 |E h-E g|\leq E|h-g|\leq2P(h\ne g).
\]
If $Eh>2P(h\ne g)$, then $Eg>0$. The factor two is sharp.
\end{proposition}
\begin{proof}
The first inequality is the triangle inequality. The second follows
pointwise because $|h-g|\leq2\mathbf1\{h\ne g\}$.
The sign claim follows from $Eg\geq Eh-E|h-g|$.
Equality in both bounds is obtained when $h=1,g=-1$ on an event of
probability $q$ and $h=g=0$ otherwise, giving a difference $2q$.
\end{proof}

This supplies an interpretable sensitivity diagnostic when a tolerance,
judge, or priority ordering changes: report the fraction of comparisons
whose decision changes and whether they reverse or become tied.
Statistical uncertainty in an estimated disagreement mass still needs
inference. Choosing a favorable hierarchy after seeing the results is
multiple selection, not prespecification. A finite family can be handled
with simultaneous confidence sequences allocated over the family, or a
separate validation stream; a result for one frozen hierarchy does not
validate arbitrary post hoc searching.

\subsection{Scope of the established guarantees}

The proofs establish identification for the stated paired design,
finite-sample error control for its bounded scores, and correctly
defined deployment conjunctions. They do not establish that an
unvalidated grader measures the intended construct; that an absence of
observed rare harms certifies safety; that a benchmark represents a
future workload; or that disjoint-pair analysis is statistically
efficient relative to a correctly analyzed all-pairs U-statistic.
Disjoint pairs intentionally simplify sequential dependence at a
possible efficiency cost. Those issues require measurement audits,
empirical comparisons, and application-specific assumptions rather than
additional algebra on the same win score.
